\begin{annexitem}{annex:p15.pe.critical}{thm:p15-pe-critical}{Constructing types and conditional critical values}
\begin{proof}
There are finitely many atomic formulas on a fixed finite variable list,
since the vocabulary and arities are finite. A type specifies their truth
values, subject to equality substitution and the equality axioms.
After collapsing its equality classes, its relation literals specify
distinct ground-atom keys on distinct class representatives. Their
probability is a product \(w_\beta>0\) of factors \(p_R\) or \(1-p_R\).
If there are \(s\) classes, the probability of assigning distinct
representatives is \((n)_s/n^s\), which tends to one; hence
\(\mu_n(\beta(\bar x_E))\to w_\beta>0\). For the original list its
equality pattern contributes \((n)_s/n^k\); conditional ratios below
cancel that factor. Thus the possibly vanishing probability of an
uncollapsed equality type is not mistaken for a positive limiting one.

Fix any assignment \(\bar a\) realizing the diagram. The probability
of \(\beta(\bar a)\) is the same \(w_\beta\) for all such assignments.
Conditional on it, parameter-only atom keys are fixed and all other
keys retain their independent Bernoulli laws. A bound tuple is fresh
if its coordinates are distinct and avoid all parameter values.
Among all \(n^\ell\) tuples the nonfresh fraction is at most
\(\ell k/n+\ell(\ell-1)/(2n)\), by summing the collision events.
For a fresh tuple, equality evaluation and identification of atom keys
are exactly those in the body construction. Finite Boolean truth
patterns partition the remaining independent bits, so their satisfying
weight is precisely the displayed finite sum \(q_\beta\).
Every nonfresh tuple has truth probability between zero and one.
The conditional mean therefore differs from \(q_\beta\) by at most
the collision fraction, uniformly in the parameter assignment.
Conditional averaging over assignments realizing \(\beta\) changes
nothing: all have the same weight and the same equality pattern.
This proves the specified conditional limit. If \(\ell=0\), the scope
is determined by the complete parameter type and its value is zero
or one, without a collision term.

The finite sum consists of products of \(p_R\) and \(1-p_R\);
expansion gives an integer-coefficient polynomial. Finitely many
types give finitely many critical values. Rational substitution gives
rational values. No independence is asserted for two tuple evaluations
which share a ground-atom key.
\end{proof}
\end{annexitem}

\begin{annexitem}{annex:p15.pe.uniform}{thm:p15-pe-uniform}{A full concentration proof for overlapping tuples}
\begin{proof}
First prove the finite bounded-difference estimate used here.
Let \(Z_1,\ldots,Z_N\) be independent finite-valued variables and let
\(F\) be a real function whose change on replacing \(Z_j\) alone is
at most \(c_j\). Successively reveal the coordinates and take finite
conditional averages: \(M_j=\mathbb E(F\mid Z_1,\ldots,Z_j)\).
For any fixed revealed past, the possible values of \(M_j-M_{j-1}\)
lie in an interval of length at most \(c_j\), and their conditional
mean is zero. Indeed the future coordinates have the same law for
each possible value of \(Z_j\); averaging preserves the change bound.

For a centered variable \(X\) in an interval of length \(c\), put
\(g(t)=\log\mathbb E e^{tX}\). Finite sums permit differentiation;
\(g''(t)\) is the variance in the exponentially tilted law.
Any law in that same interval has variance at most \(c^2/4\):
its mean minimizes squared error, and error about the midpoint is
at most \(c^2/4\). Since \(g(0)=g'(0)=0\), integration gives
\(g(t)\le t^2c^2/8\), also for negative \(t\).
Apply this conditional bound to each martingale increment and
successively average to obtain
\[
 \mathbb E e^{t(F-\mathbb EF)}
          \le \exp\bigl(t^2\textstyle\sum_jc_j^2/8\bigr).
\]
The elementary inequality
\(\Pr(W\ge a)\le\mathbb EW/a\) for \(W\ge0\), applied to the
exponential, and the choice
\(t=4\epsilon/\sum_jc_j^2\), yield
\[
 \Pr(|F-\mathbb EF|\ge\epsilon)
       \le 2\exp(-2\epsilon^2/\textstyle\sum_jc_j^2).
\]
If the sum is zero, \(F\) is constant and no tail estimate is needed.
This derivation uses only the actual finite product measure.

Fix \(\beta\) and a parameter assignment \(\bar a\) of its diagram,
and condition on \(\beta(\bar a)\). Let \(F\) be the fraction of
bound tuples satisfying \(\alpha(\bar a,\bar y)\).
Parameter-only keys are fixed. List the remaining relation-atom
patterns in the formula, after collapsing parameter equalities.
Let their number be \(h\), and let pattern \(t\) contain \(d_t\ge1\)
distinct bound variables. Replacing one ground relation bit can
affect that pattern at no more than \(n^{\ell-d_t}\) tuples, hence
can change \(F\) by at most \(n^{-d_t}\).
The pattern uses at most \(n^{d_t}\) ground keys.
If a key occurs in several patterns, sum their bounds. By
\((\sum_{t=1}^{h}b_t)^2\le h\sum_tb_t^2\), summing over keys gives
\[
 \sum_jc_j^2
   \le h\sum_t n^{d_t}n^{-2d_t}\le h^2/n.
\]
This counts repeated bound coordinates and overlapping atoms rather
than declaring their truth indicators independent. Equality atoms
are deterministic. If \(h=0\), only deterministic counting remains.

The preceding critical-value proof bounds
\(|\mathbb E(F\mid\beta(\bar a))-q_\beta|\) by a constant divided
by \(n\). Noncriticality supplies
\(\delta=\min_\beta|r-q_\beta|>0\). For sufficiently large \(n\),
deviation of \(F\) to the wrong side of \(r\) requires deviation
from its conditional mean of at least \(\delta/2\). The displayed
tail is at most \(2e^{-c n}\), with a positive constant depending
only on the fixed formula and gap. Deterministic patterns also
stay on the correct side once their collision error is below the gap.

For any assignment exactly one complete type holds. Its bad conditional
probability bounds its unconditional bad probability after multiplying
by the probability of that type. Summing over the finitely many types
and at most \(n^k\) assignments gives a constant times \(n^ke^{-cn}\),
which tends to zero (take logarithms: \(k\log n-cn\to-\infty\)).
On the complement every assignment has the displayed equivalence.
For an empty bound list the type already fixes scope truth; the result
is exact. The non-strict finite threshold and strict content comparison
are compatible precisely because the limiting tie was excluded.
\end{proof}
\end{annexitem}

\begin{annexitem}{annex:p15.pe.elimination}{thm:p15-pe-elimination}{Actual witnesses, nested binders, and finite computation}
\begin{proof}
First let \(\alpha(\bar x,y)\) be quantifier-free. Fix a type \(\beta\)
and parameter assignment, allowing repetitions. A witness equal to a
parameter value exists exactly when one of
\(\alpha(\bar x,x_i)\) is true. For fresh \(b\) outside the at most
\(k\) parameter values, simplify the scope by the type and fresh
equalities. Its success probability is \(q_\beta\), exactly, not just
asymptotically. All remaining random atoms contain \(b\). For distinct
fresh \(b,b'\) their ground-key sets are disjoint, since their arguments
use only the parameters and the respective fresh value and contain it.
Thus fresh witness trials are independent conditional on \(\beta\).
If \(q_\beta>0\), the probability of no fresh witness is at most
\((1-q_\beta)^{n-k}\le e^{-q_\beta(n-k)}\); the case \(q_\beta=1\)
has no failure once a fresh element exists.
If \(q_\beta=0\), every Boolean assignment to the remaining bits has
positive probability because all \(p_R\) are strictly between zero
and one. Consequently the simplified scope is unsatisfiable and no
fresh witness is possible in any structure of that parameter type.
Only the explicitly included parameter witnesses remain.
Finite types and \(n^k\) assignments give uniform probability tending
to one for exactly the ordinary existential content clause.

Proceed by structural induction. Quantifier-free cases are exact.
Negation preserves equivalence, and a finite conjunction fails only
if one constituent equivalence fails, so the union bound applies.
For a quantified scope, the induction event is its universally closed
equivalence, simultaneously at every parameter and bound tuple.
On that event its truth set is unchanged, hence its ordinary
witnesses and its probability counts are unchanged. Intersect with
the quantifier-free existential event just proved or with
\cref{thm:p15-pe-uniform}; their exceptional probabilities tend to zero.
This proves the quantified step and the final uniformly closed
equivalence. In particular an inner probability binder is first
replaced on all assignments; no binder exchange or independence of
its nested truth sets is used.

A quantifier-free sentence in the declared vocabulary has no relation
atoms: positive arity requires a variable and there are no constants.
It is a Boolean constant. The probability of the sentence differing
from its content tends to zero, proving the zero--one conclusion.

For effectiveness, enumerate equality partitions and consistent atomic
truth tables for the finite variable lists, substitute their values,
and compute the finite rational products and sums of
\cref{thm:p15-pe-critical}. Exact integer arithmetic decides equality
and strict order of these rationals. Traverse the finite parse tree:
reject a probability node precisely at a computed critical equality;
otherwise form its finite disjunction of types. The existential node
also computes the parameter substitutions. Truth tables reduce
Boolean output to a finite formula. The finite enumerations have
explicit bounds determined by formula length and the fixed finite
vocabulary, so loops bounded by these bounds and parse-tree recursion
are primitive recursive with standard finite encodings. This claim
does not introduce an algorithm for arbitrary real threshold names.
\end{proof}
\end{annexitem}

\begin{annexitem}{annex:p15.pe.strong}{thm:p15-pe-strong}{Constructing one class before infinitary induction}
\begin{proof}
On a fixed set of \(k\) variable symbols there are finitely many atoms,
hence finitely many Boolean truth functions of them. Choose finite
quantifier-free representatives. There are only finitely many
single-variable existential choices and finitely many choices of a
probability-bound sublist of these variables. For each representative
and each \(r\in P_0\), keep its probability elimination identity if
noncritical; keep every ordinary existential identity. Call this finite
set of universally closed identities \(F_0\). The two preceding
proofs show that each identity has model probability tending to one.
Let \(C\cap M_n\) consist of the structures satisfying all identities
in \(F_0\). A finite union bound proves \(\mu_n(C\cap M_n)\to1\).

In any one structure of \(C\), prove by well-founded induction on
formula formation that each declared noncritical formula equals its
content at every assignment. Negation is exact. For an infinitary
conjunction, every constituent identity already holds in this same
structure, so conjunction of their truth values preserves equivalence;
there is no intersection of further probabilistic events at this step.
Every quantifier-free content on this variable budget denotes a
Boolean function of the finite atomic list, and therefore has a chosen
finite representative. Substitution of equivalent scopes preserves
counts and witnesses. Its primitive identity belongs to \(F_0\);
the corresponding existential or probability step follows. An
infinitary disjunction is treated by its truth clause, or by negation.
This proves the common-class theorem for finitary and infinitary syntax.

Each individual formula in the union has its own finite budget and
threshold set and therefore one such class. This proves weak reduction;
its sentence content is a Boolean constant, giving the zero--one law.
It does not provide one common class for that entire union.

For the obstruction in the body, a real fraction \(f\) satisfies
\(f\ge1/2-1/m\) for every \(m\ge3\) exactly when \(f\ge1/2\):
if \(f<1/2\), choose \(m>1/(1/2-f)\). Thus the indicated conjunction
is exactly the half-threshold sentence. Each conjunct is noncritical,
has true content and probability tending to one. If a common class
realized all these contents, it would be contained in the
half-threshold event, whose limiting probability is \(1/2\).
It cannot have limiting probability one. The same conjunction is a
noncritical infinitary sentence with infinitely many thresholds and
limit \(1/2\), so it also witnesses failure of the unrestricted
infinitary zero--one assertion. A disjunction of the below-half
thresholds does not have this property; the printed p.~25 example
is not used in this proof.
\end{proof}
\end{annexitem}

\begin{annexitem}{annex:p15.pe.parametric}{thm:p15-pe-parametric}{Building the conditioned law and proving its elimination}
\begin{proof}
For a ground relation instance let its support be the set of elements
actually appearing in its arguments; repeated arguments do not enlarge
that set. Group all ground relation bits with exact support \(S\) into
one block. There is a uniform finite bound \(B\) on the number of bits
in a block, depending only on the vocabulary and its maximum arity.
Blocks with different supports are independent in the original model
law, since their ground-key sets are disjoint.

Consider one conjunct of \(\pi\) and a distinct assignment to its
variables, with underlying set \(S\). Every relation literal uses
all those variables, so every ground atom in this constraint belongs
to the block with support exactly \(S\). Equality literals have fixed
truth values under the distinct diagram. Group all such instances,
including every ordering of \(S\), into an allowed-configuration event
\(C_S\) in that block. Blocks not constrained by \(\pi\) have
\(C_S\) equal to their whole sample space. Thus
\[
   \{\mathcal A:\mathcal A\models\pi\}=\bigcap_S C_S.
\]
Supports beyond the finite constraint/arity bounds have no nontrivial
events. If \(C_S\) were empty, a universe large enough to contain \(S\)
would have no model of \(\pi\), contrary to the assumption. Therefore
its probability \(Z_{|S|}>0\). The notation is legitimate because
renaming elements preserves the constraint and original weights.
For a block configuration \(v\), its conditional weight is
\[
   \frac{\boldsymbol1_{C_S}(v)}
        {Z_{|S|}}\prod_{\text{bits }t\text{ in }S}
                      p_{R_t}^{v_t}(1-p_{R_t})^{1-v_t}.
\]
Finite-product factorization shows that conditioning the global
product on \(\bigcap_SC_S\) gives the product of precisely these
block laws. In particular its denominator is a positive product of
the \(Z_{|S|}\). No individual-bit independence within a constrained
block has been asserted.

A complete type on a parameter set \(A_0\) fixes all ground bits in
every block \(S\subseteq A_0\). It is consistent with \(\pi\) exactly
when these fixed configurations are allowed. Necessity follows from
the constraint. For sufficiency choose an allowed configuration in
every remaining block; these independent choices extend the type to
every larger finite universe. Such a type has positive conditional
weight
\[
  w_\beta^\pi=
       \frac{w_\beta}{\prod_{S\subseteq A_0}Z_{|S|}}>0,
\]
where empty or unconstrained blocks contribute one. Its collapsed
sampling probability tends to this weight. For a fresh bound tuple,
the joint law of all finitely many blocks on parameters and fresh
coordinates is independent of their element names and of \(n\).
Condition on the parameter type and sum the weights of allowed
configurations satisfying the scope to define \(q_\beta^\pi\).
This is an explicit finite sum of normalized products, not necessarily
the unconditioned polynomial. Nonfresh tuples have fraction \(O(1/n)\)
exactly as in \cref{thm:p15-pe-critical}, so this sum is the actual
conditional limit. The same argument proves finite critical inventories
and positive collapsed type limits. It also explains the positive
type-weight adjustment used in source Lemmas~5.2--5.4, without
confusing assignment sampling with conditioning on a fixed tuple.

For concentration condition on \(\beta(\bar a)\). The nonparameter
blocks remain independent. A relation pattern containing \(d_t\ge1\)
distinct bound variables can be affected by changing a support block
only if its ground key has that support. At most \(B\) keys in the
changed block can match this pattern, and each determines at most
\(n^{\ell-d_t}\) tuples. Thus its normalized influence is at most
\(B n^{-d_t}\); the number of support blocks it can use is at most
\(n^{d_t}\). For \(h\) patterns the same squared-sum argument gives
\(\sum_S c_S^2\le B^2h^2/n\).
The finite-valued concentration proof in
\cref{thm:p15-pe-uniform} applies to whole blocks. A positive
noncritical gap and the \(O(1/n)\) mean error give an exponential tail;
union over types and assignments gives uniform binder elimination.

For ordinary existence, the nonparameter blocks used by the scope
at a fresh single value \(b\) all contain \(b\) and otherwise only
parameter values. For different fresh \(b,b'\) these support sets
are disjoint, so the witness trials are independent under the
conditioned block law. Their exact success probability is
\(q_\beta^\pi\). A positive value supplies the exponentially small
no-witness bound; a zero value makes the scope impossible at a fresh
value because every allowed block configuration has positive weight.
Parameter witnesses are included explicitly. The type/assignment union
bound and structural induction from
\cref{thm:p15-pe-elimination} now apply, with the conditioned law and
these newly established primitive identities. This proves finitary
elimination and the sentence zero--one assertion. The construction of
the finite primitive family and its one common class from
\cref{thm:p15-pe-strong} uses only those identities and finite union
bounds, so it also proves the stated fixed-budget infinitary version.
This is a new proof of the extension, not a claim that unconditional
ground atoms stayed independent after imposing symmetry or other
parametric constraints.
\end{proof}
\end{annexitem}

\begin{annexitem}{annex:p15.pe.model-boundaries}{thm:p15-pe-model-boundaries}{Conditioned edge weights and an uncolored bipartite counterexample}
\begin{proof}
Irreflexivity is the parametric conjunct
\(\forall x\,\neg R(x,x)\); symmetry is
\(\forall x y[D(x,y)\to(R(x,y)\leftrightarrow R(y,x))]\).
Every variable occurs in each relation literal of its conjunct.
For a distinct pair, its two originally independent ordered bits have
weights \(p^2\) for \((1,1)\) and \((1-p)^2\) for \((0,0)\).
Conditioning on equality of the bits normalizes by their sum, giving
the asserted edge probability. Diagonal blocks are separately forced
to zero. Thus the global conditioning exists at every finite size
and different unordered-pair blocks remain independent.

Suppose a parametric sentence in the single binary graph vocabulary
defined exactly all uncolored finite bipartite graphs. A relation
literal can contain at most two distinct variables, and an equality
literal at most two. Hence every nonconstant constraint is on a
singleton or a pair; a Boolean constant involving a longer list is
either vacuous or excludes all sufficiently large universes, the latter
being incompatible with the supposed class.
The one-vertex loopless graph is bipartite, so its allowed singleton
configuration is the zero diagonal bit. The two-vertex empty graph
and the two-vertex one-edge undirected graph are both bipartite, so
the pair configurations \((R(a,b),R(b,a))=(0,0)\) and \((1,1)\)
are both permitted. Constraints on a pair cannot inspect other
vertices or singleton bits: every variable must occur in every
literal. Consequently the triangle with zero diagonal and every
off-diagonal bit one satisfies all the same singleton and pair
constraints. It is not bipartite: assigning one endpoint either of
two colors forces its two neighbors to the other color, although
they are adjacent. This contradicts the asserted definition.
This counterexample concerns precisely the stated uncolored
vocabulary, not an unspecified enriched presentation.

For the all-size qualification take a vocabulary with one unary
relation. Its largest relation arity is one. In
\(\forall x y[D(x,y)\to x=y]\), the equality literal contains both
variables, so it meets the parametric formation clause.
At size one its antecedent is never true, hence any unary-relation
interpretation is a model. At every size at least two, two distinct
elements make the antecedent true and the consequent false.
Thus it has a model of size equal to the largest relation arity,
but not models of all sizes. The source's p.~22 arity-only existence
criterion cannot be used with this equality-literal convention.
The selected transfer theorem assumes all positive finite sizes and
does not use that criterion to manufacture the hypothesis.

Finally, if a nullary relation is admitted as an independent ground
bit of probability \(p\in(0,1)\), the sentence consisting of that
atom has probability \(p\) at every universe size. It is already
quantifier-free and noncritical, so a zero--one assertion covering
that signature is false. The positive-arity guard avoids this exact
countermodel; it is not a supposition that asymptotic elimination
by itself forces every quantifier-free sentence to be constant.
\end{proof}
\end{annexitem}

\begin{annexitem}{annex:p15.pe.sequences}{thm:p15-pe-sequences}{Eventual margins without a hidden uniform atomic bound}
\begin{proof}
For each \(n\), fix a parameter type and any assignment of its equality
diagram. Its ground-atom probability is positive because each
\(p_R(n)\) lies in \((0,1)\), and is the same for every assignment of
that diagram. Conditional on the type, the other ground bits remain
independent. Let \(a_{\beta,n}\) be the exact fresh-tuple success
probability, computed from the finite Boolean truth-pattern sum with
\(p_R(n)\). Conditional uniform tuple averaging gives
\[
       |q_{\beta,n}-a_{\beta,n}|\le K/n
\]
for a fixed collision constant \(K\). There is no division by a
possibly small type weight in this bound: conditioning leaves parameter
assignments uniform within their equality diagram, and collision
counts are conditional counts. The mean of the truth fraction for a
fixed parameter assignment equals \(q_{\beta,n}\), by relabeling
invariance; it also differs from \(a_{\beta,n}\) by the same bound.

For each type at a noncritical probability node, failure of the two
critical inequalities means that one of
\(\liminf(q_{\beta,n}-r(n))\) and
\(\liminf(r(n)-q_{\beta,n})\) is positive. They cannot both be
positive. Consequently there are a sign, an \(\epsilon_\beta>0\)
and an \(N_\beta\) such that the appropriate signed difference is
at least \(\epsilon_\beta\) for every \(n\ge N_\beta\).
Finitely many types give one positive minimum margin and one finite
starting index. The influence bound
\(\sum_jc_j^2\le h^2/n\) from
\cref{thm:p15-pe-uniform} does not depend on the values of the
Bernoulli parameters. Its exponential-moment proof therefore gives
uniform exponential error at this node for all sufficiently large
\(n\), even if some \(p_R(n)\) tend to zero or one. The type union
bound requires no lower bound on the probability of those types.
Its chosen content is exactly the positive-liminf disjunction.

For an ordinary existential node, a type with
\(\liminf q_{\beta,n}>0\) has
\(a_{\beta,n}\ge\epsilon>0\) for all sufficiently large \(n\),
by the collision bound. The disjoint fresh witness trials give
\((1-a_{\beta,n})^{n-k}\le e^{-\epsilon(n-k)}\).
If that positive-liminf condition is not used, source
Definition~7.2 requires logical validity of the parameter-witness
disjunction under \(\beta\). Every assignment of this type already
has an actual parameter witness, in every structure; its content
is true there by that disjunction, independently of all atomic
probabilities. These two permitted cases exhaust the types at a
noncritical existential node. Thus its primitive identity is uniformly
almost-everywhere valid. There is no unsupported inference that a
merely positive but vanishing fresh-witness probability suffices.

With a finite variable budget there are finitely many Boolean
representatives, and with finitely many threshold sequences there
are finitely many primitive probability and existential identities.
Keep those meeting the stated sequence-noncritical tests. Their
eventual margins and indices can be minimized and maximized over
this finite family; their simultaneous identity class has probability
tending to one. In each structure of that class, the well-founded
infinitary induction in \cref{thm:p15-pe-strong} is purely semantic
and remains valid. This proves the common-class assertion.
For the union, choose the budget and sequence inventory of the
individual formula; weak reduction follows. Sentence contents are
Boolean constants in the positive-arity, constant-free vocabulary,
so the same comparison proves the zero--one law.

The source's comparison with numerical quantifiers is direct integer
counting: for a count \(m\),
\(m/n^\ell\ge r(n)\) iff
\(m\ge\lceil r(n)n^\ell\rceil\). In particular \(r(n)=1/n\)
and \(\ell=1\) require one witness when \(n>1\).
That semantic equality does not equate the two noncriticality tests.
It also does not extend the source's allowed threshold interval to
include the exceptional value one at \(n=1\).
\end{proof}
\end{annexitem}
